% GENERATED FILE. Edit canonical lessons, then run npm run generate:books.
% canonical-source-hash: fnv1a64:fe62eea719d0dd69
% canonical-lessons: JA-C373-kaisatsu, JA-C373-teikiken, JA-C373-kaisuuken, JA-C373-seisan, JA-C373-tourisugiru

\chapter{Through the Ticket Barrier}
\label{ch:ja-through-the-ticket-barrier}
\hlchaptermodality{373}

\begin{chapteropening}
\textbf{By the end of this chapter:} \emph{I can get through the ticket barrier and pay the right fare.}

Complete the last lesson of chapter 373: I can get through the ticket barrier and pay the right fare.
\end{chapteropening}

\section[\texorpdfstring{\ja{かいさつ} (kaisatsu)}{かいさつ (kaisatsu)}]{\ja{かいさつ} (kaisatsu) — ticket inspection; the ticket barrier}

\label{lesson:JA-C373-kaisatsu}



\begin{quote}
Before the new one: say the Japanese for a waiting room, then the Japanese for a vending machine.
\end{quote}

\subsection*{You'll want to know: \ja{かいさつ}}
\textbf{\ja{かいさつ}} — \emph{kaisatsu} — \textquotedblleft{}ticket inspection; the ticket barrier\textquotedblright{}.

\begin{quote}
\textbf{\ja{かいさつの} \ja{まえで} \ja{あにを} \ja{まちました。}} — \emph{kaisatsu no mae de ani o machimashita} — \textquotedblleft{}I waited for my brother in front of the ticket barrier.\textquotedblright{}
\end{quote}

\begin{cousinweb}
\textbf{\ja{かいさつ}} (\emph{kaisatsu}) is Sino-Japanese (kango): two Chinese characters read with their on'yomi, readings borrowed from Chinese. They are \emph{kai}, \textquotedblleft{}examine\textquotedblright{} (Middle Chinese \emph{kojX}) and \emph{satsu}, \textquotedblleft{}tag\textquotedblright{} (\emph{tsreat}) (KANJIDIC2; Middle Chinese in Baxter's transcription).
\end{cousinweb}

\begin{grammarlens}[title={What you've built}]
One more word for getting around.
\end{grammarlens}

\subsection*{Your turn}
\begin{itemize}
\raggedright
  \item \emph{Say it:} \emph{kaisatsu}
  \item \emph{Say it:} \emph{kaisatsu}, once more
  \item \emph{Say it:} \emph{kaisatsu no mae de ani o machimashita} — I waited for my brother in front of the ticket barrier.
  \item \emph{Recall:} say the Japanese for a waiting room, then the Japanese for a vending machine, then say \emph{kaisatsu} again
\end{itemize}

\begin{tcolorbox}[breakable,colback=teal!4,colframe=teal!35!black,title={Before you move on}]
What does \ja{かいさつ} mean? (\textquotedblleft{}Ticket inspection; the ticket barrier\textquotedblright{}.) Say it once more.
\end{tcolorbox}
\section[\texorpdfstring{\ja{ていきけん} (teikiken)}{ていきけん (teikiken)}]{\ja{ていきけん} (teikiken) — a commuter pass}

\label{lesson:JA-C373-teikiken}



\begin{quote}
Before the new one: say the Japanese for a vending machine, then the Japanese for ticket inspection.
\end{quote}

\subsection*{You'll want to know: \ja{ていきけん}}
\textbf{\ja{ていきけん}} — \emph{teikiken} — \textquotedblleft{}a commuter pass\textquotedblright{}.

\begin{quote}
\textbf{\ja{ちちは} \ja{ていきけんで} \ja{でんしゃに} \ja{のります。}} — \emph{chichi wa teikiken de densha ni norimasu} — \textquotedblleft{}My father rides the train on a commuter pass.\textquotedblright{}
\end{quote}

\begin{cousinweb}
\textbf{\ja{ていきけん}} (\emph{teikiken}) is Sino-Japanese (kango): three Chinese characters read with their on'yomi, readings borrowed from Chinese. They are \emph{tei}, \textquotedblleft{}fix\textquotedblright{} (Middle Chinese \emph{dengH}), \emph{ki}, \textquotedblleft{}period\textquotedblright{} (\emph{gi}) and \emph{ken}, \textquotedblleft{}ticket\textquotedblright{} (\emph{khjwonH}) (KANJIDIC2; Middle Chinese in Baxter's transcription).
\end{cousinweb}

\begin{grammarlens}[title={What you've built}]
One more word for getting around.
\end{grammarlens}

\subsection*{Your turn}
\begin{itemize}
\raggedright
  \item \emph{Say it:} \emph{teikiken}
  \item \emph{Say it:} \emph{teikiken}, once more
  \item \emph{Say it:} \emph{chichi wa teikiken de densha ni norimasu} — My father rides the train on a commuter pass.
  \item \emph{Recall:} say the Japanese for a vending machine, then the Japanese for ticket inspection, then say \emph{teikiken} again
\end{itemize}

\begin{tcolorbox}[breakable,colback=teal!4,colframe=teal!35!black,title={Before you move on}]
What does \ja{ていきけん} mean? (\textquotedblleft{}A commuter pass\textquotedblright{}.) Say it once more.
\end{tcolorbox}
\section[\texorpdfstring{\ja{かいすうけん} (kaisūken)}{かいすうけん (kaisūken)}]{\ja{かいすうけん} (kaisūken) — a book of tickets}

\label{lesson:JA-C373-kaisuuken}



\begin{quote}
Before the new one: say the Japanese for ticket inspection, then the Japanese for a commuter pass.
\end{quote}

\subsection*{You'll want to know: \ja{かいすうけん}}
\textbf{\ja{かいすうけん}} — \emph{kaisūken} — \textquotedblleft{}a book of tickets\textquotedblright{}.

\begin{quote}
\textbf{\ja{かいすうけんは} \ja{ふつうの} \ja{きっぷより} \ja{やすいです。}} — \emph{kaisūken wa futsū no kippu yori yasui desu} — \textquotedblleft{}A book of tickets is cheaper than ordinary tickets.\textquotedblright{}
\end{quote}

\begin{cousinweb}
\textbf{\ja{かいすうけん}} (\emph{kaisūken}) is Sino-Japanese (kango): three Chinese characters read with their on'yomi, readings borrowed from Chinese. They are \emph{kai}, \textquotedblleft{}round\textquotedblright{} (Middle Chinese \emph{hwoj}), \emph{sū}, \textquotedblleft{}number\textquotedblright{} (\emph{srjuH}) and \emph{ken}, \textquotedblleft{}ticket\textquotedblright{} (\emph{khjwonH}) (KANJIDIC2; Middle Chinese in Baxter's transcription).
\end{cousinweb}

\begin{grammarlens}[title={What you've built}]
One more word for getting around.
\end{grammarlens}

\subsection*{Your turn}
\begin{itemize}
\raggedright
  \item \emph{Say it:} \emph{kaisūken}
  \item \emph{Say it:} \emph{kaisūken}, once more
  \item \emph{Say it:} \emph{kaisūken wa futsū no kippu yori yasui desu} — A book of tickets is cheaper than ordinary tickets.
  \item \emph{Recall:} say the Japanese for ticket inspection, then the Japanese for a commuter pass, then say \emph{kaisūken} again
\end{itemize}

\begin{tcolorbox}[breakable,colback=teal!4,colframe=teal!35!black,title={Before you move on}]
What does \ja{かいすうけん} mean? (\textquotedblleft{}A book of tickets\textquotedblright{}.) Say it once more.
\end{tcolorbox}
\section[\texorpdfstring{\ja{せいさん} (seisan)}{せいさん (seisan)}]{\ja{せいさん} (seisan) — fare adjustment; settling up}

\label{lesson:JA-C373-seisan}



\begin{quote}
Before the new one: say the Japanese for a commuter pass, then the Japanese for a book of tickets.
\end{quote}

\subsection*{You'll want to know: \ja{せいさん}}
\textbf{\ja{せいさん}} — \emph{seisan} — \textquotedblleft{}fare adjustment; settling up\textquotedblright{}.

\begin{quote}
\textbf{\ja{えきの} \ja{まどぐちで} \ja{せいさんを} \ja{しました。}} — \emph{eki no madoguchi de seisan o shimashita} — \textquotedblleft{}I paid the fare difference at the station window.\textquotedblright{}
\end{quote}

\begin{cousinweb}
\textbf{\ja{せいさん}} (\emph{seisan}) is Sino-Japanese (kango): two Chinese characters read with their on'yomi, readings borrowed from Chinese. They are \emph{sei}, \textquotedblleft{}refined\textquotedblright{} (Middle Chinese \emph{tsjeng}) and \emph{san}, \textquotedblleft{}calculate\textquotedblright{} (\emph{swanX}) (KANJIDIC2; Middle Chinese in Baxter's transcription).
\end{cousinweb}

\begin{grammarlens}[title={What you've built}]
One more word for getting around.
\end{grammarlens}

\subsection*{Your turn}
\begin{itemize}
\raggedright
  \item \emph{Say it:} \emph{seisan}
  \item \emph{Say it:} \emph{seisan}, once more
  \item \emph{Say it:} \emph{eki no madoguchi de seisan o shimashita} — I paid the fare difference at the station window.
  \item \emph{Recall:} say the Japanese for a commuter pass, then the Japanese for a book of tickets, then say \emph{seisan} again
\end{itemize}

\begin{tcolorbox}[breakable,colback=teal!4,colframe=teal!35!black,title={Before you move on}]
What does \ja{せいさん} mean? (\textquotedblleft{}Fare adjustment; settling up\textquotedblright{}.) Say it once more.
\end{tcolorbox}
\section[\texorpdfstring{\ja{とおりすぎる} (tōrisugiru)}{とおりすぎる (tōrisugiru)}]{\ja{とおりすぎる} (tōrisugiru) — to pass by; to go past}

\label{lesson:JA-C373-tourisugiru}



\begin{quote}
Before the new one: say the Japanese for a book of tickets, then the Japanese for fare adjustment.
\end{quote}

\subsection*{You'll want to know: \ja{とおりすぎる}}
\textbf{\ja{とおりすぎる}} — \emph{tōrisugiru} — \textquotedblleft{}to pass by; to go past\textquotedblright{}.

\begin{quote}
\textbf{\ja{でんしゃは} \ja{ちいさい} \ja{えきを} \ja{とおりすぎました。}} — \emph{densha wa chiisai eki o tōrisugimashita} — \textquotedblleft{}The train went past the small station.\textquotedblright{}
\end{quote}

Polite \textbf{\ja{とおりすぎます}}, \emph{tōrisugimasu}; an ichidan verb. Past (polite): \textbf{\ja{とっきゅうが} \ja{えきを} \ja{とおりすぎました。}}, \emph{tokkyū ga eki o tōrisugimashita}.

\begin{cousinweb}
\textbf{\ja{とおりすぎる}} (\emph{tōrisugiru}) is a native Japanese word (wago), inherited rather than borrowed from Chinese, built from \emph{tōri}, \textquotedblleft{}passing\textquotedblright{} (from \emph{tōru}) and \emph{sugiru}, \textquotedblleft{}to go beyond\textquotedblright{}; its kanji were attached later as kun'yomi (JMdict).
\end{cousinweb}

\begin{grammarlens}[title={What you've built}]
One more word for getting around.
\end{grammarlens}

\subsection*{Your turn}
\begin{itemize}
\raggedright
  \item \emph{Say it:} \emph{tōrisugiru}
  \item \emph{Say it:} \emph{tōrisugiru}, once more
  \item \emph{Say it:} \emph{densha wa chiisai eki o tōrisugimashita} — The train went past the small station.
  \item \emph{Recall:} say the Japanese for a book of tickets, then the Japanese for fare adjustment, then say \emph{tōrisugiru} again
\end{itemize}

\begin{tcolorbox}[breakable,colback=teal!4,colframe=teal!35!black,title={Before you move on}]
What does \ja{とおりすぎる} mean? (\textquotedblleft{}To pass by; to go past\textquotedblright{}.) Say it once more.
\end{tcolorbox}
